% ================================================================
% EXERCICE INDIVIDUEL — ECNSA
% Structure calquée sur exercice_individuel.tex du Chap2.
%
% UTILISATION :
%   1. Copier ce fichier dans le dossier du devoir.
%   2. Remplacer TOUS les champs [entre crochets].
%   3. Compiler avec pdflatex (une seule passe suffit).
%
% LIMITE : ~5 pages. Pas de fancyhdr, numéro de page en bas centré.
% Police 10pt et marges compactes pour maximiser l'espace disponible.
% ================================================================
\documentclass[a4paper,10pt]{article}

% Chargement du préambule — cherche le fichier dans plusieurs chemins
% relatifs pour fonctionner depuis n'importe quel sous-dossier du projet.
\newcommand{\chargerpreambule}[1]{%
  \IfFileExists{../../templates/#1}{\input{../../templates/#1}}{%
  \IfFileExists{../templates/#1}{\input{../templates/#1}}{%
  \IfFileExists{templates/#1}{\input{templates/#1}}{\input{#1}}}}}
\chargerpreambule{preamble-commun.tex}
\chargerpreambule{preamble-exercice-individuel.tex}

% ----------------------------------------------------------------
% MÉTADONNÉES — à personnaliser
% ----------------------------------------------------------------
\title{\textbf{Exercice individuel : [titre]}}
\author{[Prénom NOM]\\ Classe : [classe]}
\date{[date]}     % ou \date{\today} pour la date du jour

\begin{document}
  \maketitle

  % ================================================================
  % SECTION 1 — Objet et périmètre (toujours présente)
  % ================================================================
  \section{Objet et périmètre}

  [Présenter la consigne, le contexte réglementaire ou technique, et le périmètre
  de l'analyse. Préciser les hypothèses retenues si nécessaire.]

  % ================================================================
  % SECTION 2 — Premier thème : obligations et écarts
  % Renommer le titre selon l'énoncé.
  % ================================================================
  \section{[Thème 1 — ex. Obligations applicables et écarts constatés]}

  % Dupliquer les blocs \subsection pour chaque obligation à examiner.

  \subsection{[Sous-thème 1]}

  \textbf{Obligation.} [Énoncer la règle ou l'obligation applicable.]

  \textbf{Constat.} [Décrire la situation observée et l'écart par rapport
  à l'obligation. Rester factuel.]

  \subsection{[Sous-thème 2]}

  \textbf{Obligation.} [Énoncer la règle ou l'obligation applicable.]

  \textbf{Constat.} [Décrire la situation observée et l'écart.]

  \subsection{[Sous-thème 3]}

  \textbf{Obligation.} [Énoncer la règle ou l'obligation applicable.]

  \textbf{Constat.} [Décrire la situation observée et l'écart.]

  % ================================================================
  % SECTION 3 — Deuxième thème : mesures de mise en conformité
  % ================================================================
  \section{[Thème 2 — ex. Mesures de mise en conformité]}

  % Dupliquer les blocs \subsection pour chaque mesure supplémentaire.

  \subsection{Mesure 1 : [intitulé]}

  \textbf{Règle.} [Énoncer la règle ou l'exigence à respecter.]

  \textbf{Mise en œuvre.} [Décrire les actions concrètes et les acteurs responsables.]

  \textbf{Preuve.} [Indiquer les documents, journaux ou résultats permettant
  de vérifier que la mesure est appliquée.]

  \subsection{Mesure 2 : [intitulé]}

  \textbf{Règle.} [Règle applicable.]

  \textbf{Mise en œuvre.} [Actions concrètes et responsables.]

  \textbf{Preuve.} [Documents ou résultats de vérification.]

  % ================================================================
  % SECTION 4 — Tableau de synthèse (toujours présent en conclusion)
  % Une ligne par obligation / mesure identifiée ci-dessus.
  % ================================================================
  \section{Responsable et preuve de la mise en œuvre}

  \begin{table}[h]
    \centering
    \small
    \begin{tabularx}{\textwidth}{@{}XXXX@{}}
      \toprule
      \textbf{Obligation} & \textbf{Mesure} & \textbf{Responsable} & \textbf{Preuve} \\
      \midrule
      [Obligation 1] & M1 : [intitulé court] & [Rôle(s)] & [Preuve attendue] \\
      [Obligation 2] & M2 : [intitulé court] & [Rôle(s)] & [Preuve attendue] \\
      [Obligation 3] & M3 : [intitulé court] & [Rôle(s)] & [Preuve attendue] \\
      % Ajouter une ligne par obligation supplémentaire.
      \bottomrule
    \end{tabularx}
  \end{table}

\end{document}
